\section{Limitations and ethics}

This is a pilot benchmark, not a representative sample of multilingual or
code-switched use. It covers only three language pairs and uses synthetic
prompts. The English content embedded in the items reflects a common editing
workflow, but many global workflows involve non-English target content,
non-Latin scripts, dialectal variation, or richer social-register distinctions
that are not covered here.

The automatic scorer is intentionally conservative and strongest for exact
span-preservation and script checks. It is weaker for nuanced register and
cultural appropriateness. The blinded LLM-judge audit supports the sampled
pass/fail labels, but it is not a substitute for native/near-native human
validation. Any release or final workshop version should clearly label the
human audit status and avoid claims that the benchmark validates cultural
appropriateness for all included communities.

The WildChat cue scan uses only aggregate counts and hashed metadata. We do not
release raw user logs or quote private conversation text. The scan is useful for
motivating the taxonomy, but regex cue hits are not reliable prevalence
estimates. The benchmark items are synthetic and should be released with a
benchmark card explaining intended use, non-representativeness, language-pair
coverage, and known scorer limitations.

Finally, token tax is a normalized repair-burden metric, not a direct economic
claim. In our experiments, the longer contract prompt reduces repair-normalized
token tax but increases absolute token use. Reporting both numbers is important
for users whose constraints are monetary cost, latency, or context-window
budget.
